% ============================================================
\chapter{Model Deployment --- REST APIs, FastAPI, Streamlit}
\label{ch:deployment}
\index{deployment}
% ============================================================

\section{Deployment Patterns}
\label{sec:deployment-patterns}

Once a model is trained, it must be made available for inference.
Three major patterns exist:

\begin{definition}[Deployment pattern]
A \textbf{deployment pattern}\index{deployment!pattern} defines how a
trained model serves predictions: in batch on stored data, in real time
via an API, or continuously on a data stream.
\end{definition}

\begin{center}
\begin{tabular}{@{}llp{5.5cm}l@{}}
\toprule
\textbf{Pattern} & \textbf{Latency} & \textbf{Use Case} & \textbf{Example} \\
\midrule
Batch & Minutes--hours & Nightly scoring, report generation &
  Spark job, Airflow \\
Online (REST) & Milliseconds & Real-time predictions, user-facing
  applications & FastAPI, TorchServe \\
Streaming & Sub-second & Fraud detection, real-time
  recommendations & Kafka, Flink \\
\bottomrule
\end{tabular}
\end{center}

\begin{center}
\begin{tikzpicture}[
  box/.style={rectangle, draw, rounded corners, minimum width=2.5cm,
    minimum height=1.2cm, align=center, font=\small},
  arr/.style={-{Stealth[length=2.5mm]}, thick, steelblue}
]
  % Client
  \node[box, fill=blue!10] (client) {\textbf{Client}\\Web / Mobile};
  % API Gateway
  \node[box, fill=cyan!10, right=2cm of client] (gw) {\textbf{API}\\
    \textbf{Gateway}};
  % Model service
  \node[box, fill=green!10, right=2cm of gw] (model) {\textbf{Model}\\
    \textbf{Service}};
  % Model store
  \node[box, fill=orange!10, above=1cm of model] (store) {\textbf{Model}\\
    \textbf{Store}};
  % Feature store
  \node[box, fill=purple!10, below=1cm of model] (feat) {\textbf{Feature}\\
    \textbf{Store}};
  % Monitoring
  \node[box, fill=red!10, right=2cm of model] (mon) {\textbf{Monitoring}\\
    Prometheus};

  \draw[arr] (client) -- node[above, font=\scriptsize] {REST} (gw);
  \draw[arr] (gw) -- node[above, font=\scriptsize] {JSON} (model);
  \draw[arr] (store) -- node[right, font=\scriptsize] {load} (model);
  \draw[arr] (feat) -- node[right, font=\scriptsize] {features} (model);
  \draw[arr] (model) -- node[above, font=\scriptsize] {metrics} (mon);
\end{tikzpicture}
\end{center}

\section{REST API Fundamentals for ML}
\label{sec:rest-api}
\index{REST API}

A REST API exposes the model through HTTP endpoints. The typical
contract for an ML service:

\begin{center}
\begin{tabular}{@{}lllp{4.5cm}@{}}
\toprule
\textbf{Endpoint} & \textbf{Method} & \textbf{Status} & \textbf{Description} \\
\midrule
\texttt{/health} & GET & 200 & Service health check \\
\texttt{/predict} & POST & 200 & Single prediction \\
\texttt{/predict/batch} & POST & 200 & Batch prediction \\
\texttt{/model/info} & GET & 200 & Model metadata (version, metrics) \\
\bottomrule
\end{tabular}
\end{center}

\section{Model Serialization}
\label{sec:serialization}
\index{serialization}

Before serving, the model must be serialized to a portable format.

\begin{center}
\begin{tabular}{@{}lllp{4cm}@{}}
\toprule
\textbf{Format} & \textbf{Framework} & \textbf{Speed} &
  \textbf{Notes} \\
\midrule
joblib / pickle & scikit-learn & Fast & Python-only, security risk \\
ONNX & Any $\to$ ONNX Runtime & Very fast & Cross-framework,
  optimized \\
TorchScript & PyTorch & Fast & No Python dependency at inference \\
SavedModel & TensorFlow & Fast & TF Serving compatible \\
\bottomrule
\end{tabular}
\end{center}

\begin{codeblock}[title=\textbf{Model serialization examples}]
\begin{minted}{python}
import joblib
import torch
import onnxruntime as ort

# --- scikit-learn with joblib ---
joblib.dump(model, "model.joblib")
model = joblib.load("model.joblib")

# --- PyTorch with TorchScript ---
scripted = torch.jit.script(model)
scripted.save("model.pt")
loaded = torch.jit.load("model.pt")

# --- ONNX export from PyTorch ---
dummy_input = torch.randn(1, 10)
torch.onnx.export(model, dummy_input, "model.onnx",
                   input_names=["input"],
                   output_names=["output"],
                   dynamic_axes={"input": {0: "batch"}})

# --- ONNX inference ---
session = ort.InferenceSession("model.onnx")
result = session.run(None, {"input": input_array})
\end{minted}
\end{codeblock}

\section{FastAPI for ML Serving}
\label{sec:fastapi}
\index{FastAPI}

\textbf{FastAPI}\index{FastAPI} is a modern, high-performance Python
web framework ideal for ML serving: automatic validation via Pydantic,
async support, and auto-generated OpenAPI documentation.

\begin{codeblock}[title=\textbf{Complete FastAPI ML serving application}]
\begin{minted}{python}
"""ML model serving API with FastAPI."""
import time
from contextlib import asynccontextmanager
from pathlib import Path

import joblib
import numpy as np
from fastapi import FastAPI, HTTPException
from pydantic import BaseModel, Field

# ---------- Pydantic schemas ----------

class PredictionRequest(BaseModel):
    features: list[float] = Field(
        ..., min_length=1,
        description="Input feature vector"
    )

class PredictionResponse(BaseModel):
    prediction: int
    probability: float
    model_version: str
    latency_ms: float

class HealthResponse(BaseModel):
    status: str
    model_loaded: bool
    model_version: str

class BatchRequest(BaseModel):
    instances: list[list[float]] = Field(
        ..., min_length=1,
        description="List of feature vectors"
    )

class BatchResponse(BaseModel):
    predictions: list[int]
    probabilities: list[float]
    count: int

# ---------- Application ----------

MODEL_PATH = Path("models/model.joblib")
MODEL_VERSION = "1.0.0"
ml_model = None

@asynccontextmanager
async def lifespan(app: FastAPI):
    """Load model on startup, cleanup on shutdown."""
    global ml_model
    if not MODEL_PATH.exists():
        raise FileNotFoundError(f"Model not found: {MODEL_PATH}")
    ml_model = joblib.load(MODEL_PATH)
    print(f"Model loaded from {MODEL_PATH}")
    yield
    ml_model = None

app = FastAPI(
    title="ML Prediction API",
    version=MODEL_VERSION,
    lifespan=lifespan,
)

@app.get("/health", response_model=HealthResponse)
async def health():
    """Health check endpoint."""
    return HealthResponse(
        status="healthy",
        model_loaded=ml_model is not None,
        model_version=MODEL_VERSION,
    )

@app.post("/predict", response_model=PredictionResponse)
async def predict(request: PredictionRequest):
    """Single prediction endpoint."""
    if ml_model is None:
        raise HTTPException(503, "Model not loaded")
    start = time.perf_counter()
    X = np.array(request.features).reshape(1, -1)
    pred = int(ml_model.predict(X)[0])
    proba = float(ml_model.predict_proba(X).max())
    latency = (time.perf_counter() - start) * 1000
    return PredictionResponse(
        prediction=pred,
        probability=proba,
        model_version=MODEL_VERSION,
        latency_ms=round(latency, 2),
    )

@app.post("/predict/batch", response_model=BatchResponse)
async def predict_batch(request: BatchRequest):
    """Batch prediction endpoint."""
    if ml_model is None:
        raise HTTPException(503, "Model not loaded")
    X = np.array(request.instances)
    preds = ml_model.predict(X).tolist()
    probas = ml_model.predict_proba(X).max(axis=1).tolist()
    return BatchResponse(
        predictions=preds,
        probabilities=[round(p, 4) for p in probas],
        count=len(preds),
    )
\end{minted}
\end{codeblock}

\begin{codeblock}[title=\textbf{Running and testing the API}]
\begin{minted}{bash}
# Run the API
uvicorn src.api:app --host 0.0.0.0 --port 8000 --reload

# Test health endpoint
curl http://localhost:8000/health

# Test prediction
curl -X POST http://localhost:8000/predict \
  -H "Content-Type: application/json" \
  -d '{"features": [5.1, 3.5, 1.4, 0.2]}'

# Interactive docs
# Open http://localhost:8000/docs in a browser
\end{minted}
\end{codeblock}

\section{Streamlit for ML Demos}
\label{sec:streamlit}
\index{Streamlit}

\textbf{Streamlit}\index{Streamlit} enables rapid creation of
interactive ML demos without front-end knowledge.

\begin{codeblock}[title=\textbf{Streamlit ML demo application}]
\begin{minted}{python}
"""Interactive ML demo with Streamlit."""
import streamlit as st
import requests
import pandas as pd

API_URL = "http://localhost:8000"

st.title("ML Model Explorer")
st.markdown("Interactive prediction demo")

# Sidebar for input
st.sidebar.header("Input Features")
feature_1 = st.sidebar.slider("Sepal Length", 4.0, 8.0, 5.1)
feature_2 = st.sidebar.slider("Sepal Width", 2.0, 4.5, 3.5)
feature_3 = st.sidebar.slider("Petal Length", 1.0, 7.0, 1.4)
feature_4 = st.sidebar.slider("Petal Width", 0.1, 2.5, 0.2)

if st.button("Predict"):
    response = requests.post(
        f"{API_URL}/predict",
        json={"features": [feature_1, feature_2,
                           feature_3, feature_4]},
    )
    if response.status_code == 200:
        result = response.json()
        st.success(f"Prediction: {result['prediction']}")
        st.metric("Confidence", f"{result['probability']:.2%}")
        st.metric("Latency", f"{result['latency_ms']:.1f} ms")
    else:
        st.error(f"API error: {response.status_code}")
\end{minted}
\end{codeblock}

\section{Serving Framework Comparison}
\label{sec:serving-comparison}

\begin{center}
\begin{tabular}{@{}lcccc@{}}
\toprule
\textbf{Criterion} & \textbf{FastAPI} & \textbf{Flask} &
  \textbf{TorchServe} & \textbf{Triton} \\
\midrule
Ease of use & High & High & Medium & Low \\
Performance & High & Medium & High & Very high \\
Auto-scaling & Manual & Manual & Built-in & Built-in \\
Multi-model & Manual & Manual & Yes & Yes \\
GPU batching & Manual & No & Yes & Yes \\
Framework lock-in & None & None & PyTorch & None \\
OpenAPI docs & Auto & Manual & No & No \\
Best for & Custom APIs & Simple APIs & PyTorch prod & High-perf prod \\
\bottomrule
\end{tabular}
\end{center}

\begin{remark}
FastAPI is the best starting point for most ML projects: it is simple,
performant, and flexible. Migrate to TorchServe or Triton when you
need advanced features like dynamic batching, multi-model serving, or
GPU-optimized inference at scale.
\end{remark}

\section{Best Practices and Anti-patterns}

\begin{bestpractice}[title=\textbf{Deployment best practices}]
\begin{enumerate}
  \item \textbf{Version your models}: include the model version in
    every response.
  \item \textbf{Validate inputs}: use Pydantic schemas with
    constraints.
  \item \textbf{Health checks}: always expose a \texttt{/health}
    endpoint.
  \item \textbf{Measure latency}: return inference time in the
    response.
  \item \textbf{Graceful startup}: load the model before accepting
    requests (lifespan events).
  \item \textbf{Log predictions}: store inputs and outputs for
    monitoring and debugging.
  \item \textbf{Separate model and API}: the API code should not
    contain training logic.
\end{enumerate}
\end{bestpractice}

\begin{antipattern}[title=\textbf{Deployment anti-patterns}]
\begin{itemize}
  \item Loading the model on every request instead of at startup.
  \item No input validation---malformed data causes cryptic errors.
  \item Returning raw numpy arrays instead of JSON-serializable types.
  \item Hardcoding model paths without environment variable overrides.
  \item No error handling---a failed prediction crashes the server.
  \item Running the API without HTTPS in production.
\end{itemize}
\end{antipattern}

\section{Mini-project: End-to-end Deployment}
\label{sec:mini-project-ch8}

\begin{miniproject}[title=\textbf{Mini-project 8: Deploy an ML model}]
\begin{enumerate}
  \item Train a scikit-learn classifier and serialize it with
    \pkg{joblib}.
  \item Build a FastAPI application with \texttt{/health},
    \texttt{/predict}, and \texttt{/predict/batch} endpoints.
  \item Add Pydantic input validation with meaningful error messages.
  \item Write a Streamlit front-end that calls the API.
  \item Containerize both services with Docker Compose.
  \item Export the model to ONNX format and compare inference latency
    against joblib.
\end{enumerate}
\end{miniproject}

\section{Exercises}
\label{sec:exercises-ch8}

\begin{exercise}[$\bigstar$ --- Basic API]
Create a FastAPI application that serves a pre-trained
\pkg{scikit-learn} model. The API must have a \texttt{/health}
endpoint and a \texttt{/predict} endpoint with Pydantic validation.
Test with \cli{curl} or the automatic Swagger documentation.
\end{exercise}

\begin{exercise}[$\bigstar\bigstar$ --- ONNX conversion]
Take a trained PyTorch model and:
\begin{enumerate}
  \item Export it to ONNX format with dynamic batch size
  \item Build a FastAPI endpoint that uses ONNX Runtime for inference
  \item Benchmark the latency of PyTorch vs ONNX Runtime on 100
    requests
  \item Document the performance difference
\end{enumerate}
\end{exercise}

\begin{exercise}[$\bigstar\bigstar\bigstar$ --- Production-grade API]
Build a production-ready ML API with:
\begin{enumerate}
  \item FastAPI with async endpoints
  \item A/B testing support (serve two model versions, route by header)
  \item Request logging to a database (async writes)
  \item Rate limiting with \pkg{slowapi}
  \item Prometheus metrics endpoint (\texttt{/metrics})
  \item Docker Compose deployment with Nginx reverse proxy
\end{enumerate}
\end{exercise}

\section{Cheatsheet}

\begin{cheatsheet}[title=\textbf{Chapter 8 Summary}]
\begin{tabular}{@{}p{4cm}p{8.5cm}@{}}
\toprule
\textbf{Concept} & \textbf{Key Point} \\
\midrule
Batch vs Online & Batch = scheduled; Online = real-time REST API \\
FastAPI & High-perf Python API with Pydantic + async \\
Pydantic & Input/output validation via type annotations \\
Model formats & joblib (sklearn), ONNX (cross-framework), TorchScript \\
Streamlit & Rapid ML demos with pure Python \\
Health check & Always expose \texttt{GET /health} \\
Lifespan events & Load model at startup, not per request \\
\bottomrule
\end{tabular}
\end{cheatsheet}
